% Lösungen Einheit 3
\einheitenkopf{Einheit 3 von 4 · Sachaufgaben}
\erg{22}{a) $x_1$ \quad b) $x_2$ \quad c) beide \quad d) $x_1$ \quad e) beide}
\erg{23}{a) $x_1 = 15$, $x_2 = -15$ \quad b) $x_1 = 30$, $x_2 = -30$ \quad c) $x^2 = 25$; $x_1 = 5$, $x_2 = -5$ \quad d) $x^2 = 64$; $x_1 = 8$, $x_2 = -8$ \quad e) $x^2 - 19 = 81$, $x^2 = 100$; $x_1 = 10$, $x_2 = -10$ \quad f) $x \cdot (x + 1) = 72$, $x^2 + x - 72 = 0$; $-0{,}5 \pm 8{,}5$; $x_1 = 8$, $x_2 = -9$; Zahlenpaare $8$ und $9$ oder $-9$ und $-8$}
\erg{24}{a) $x \cdot (x + 2) = 48$ \quad b) $x^2 + 2x - 48 = 0$; $-1 \pm 7$; $x_1 = 6$, $x_2 = -8$ \quad c) $x_1 = 6$; eine Länge ist nicht negativ \quad d) Das Rechteck ist $6\,$cm breit und $8\,$cm lang. \quad e) $6\,\text{cm} \cdot 8\,\text{cm} = 48\,\text{cm}^2$, und $8\,$cm ist $2\,$cm mehr als $6\,$cm.}
\erg{25}{a) $(x + 6)^2 = 625$; $x + 6 = 25$ oder $x + 6 = -25$; $x_1 = 19$, $x_2 = -31$; Seitenlänge $19\,$cm \quad b) $(x + 1) \cdot (x + 6) = 84$, $x^2 + 7x - 78 = 0$; $-3{,}5 \pm 9{,}5$; $x_1 = 6$, $x_2 = -13$; das Beet war $6\,$m breit und $11\,$m lang.}
\erg{26}{a) Eine Länge kann nicht negativ sein, $-9$ entfällt: Das Rechteck ist $5\,$m breit und $9\,$m lang. \quad b) $x^2 + 5x - 14 = 0$; $x_1 = 2$, $x_2 = -7$; $2\,$m breit, $7\,$m lang \quad c) Paul hat den Umfang statt des Flächeninhalts benutzt; richtig: $x \cdot (x + 3) = 54$, $x^2 + 3x - 54 = 0$; $x_1 = 6$, $x_2 = -9$; $6\,$cm breit, $9\,$cm lang \quad d) $x^2 + x - 30 = 0$; $x_1 = 5$, $x_2 = -6$; $5\,$m breit, $6\,$m lang}
\erg{27}{a) Beim Zahlenrätsel sind auch negative Zahlen erlaubt: $4^2 + 4 \cdot 4 = 32$ und $(-8)^2 + 4 \cdot (-8) = 32$. \quad b) Eine Seitenlänge ist nie negativ, $-8$ entfällt: $4\,$cm breit und $8\,$cm lang.}
